\scanpage{37}
\mspage{33}
% archive sheet number: 36
% source lines 1--10
Elle permet, en plus des \deleted{l'atelier des} figures compactes (et des figures finies dont on disposait déjà avant), d'introduire les figures localement finies. Il \unsure{ne|me} semble pas qu'en pratique, il y ait lieu de travailler avec d'autres ateliers attenants que celui des figures finies, celui des figures \uline{localement finies}, ou (à la rigueur) celui de \uline{toutes} les $M$-figures.

% source lines 11--14
\uline{NB} Les \uline{figures compactes} forment un atelier attenant à $M_c$, qui est un sous-\uline{\unsure{atelier|magasin}} de $M$, mais non à $M$, sauf si $M_c\unsure{=|\equiv}M$.

% source lines 15--19
Ainsi, on peut dire que la base \inserted{technique} de l'approche combinatoire de la topologie que je propose, est plutôt \deleted{la} notion de \uline{magasin compactologique}
\[
(M,\trianglelefteq,\overset{\circ}{\ll},\mathrel{\mid\circ\mid},M_c).
\]
% source lines 20--22
C'est une donnée plus simple que celle d'un atelier attenant à $M$, car $M_c\in\unsure{\mathcal P|\mathscr P}(M)$, et $\unsure{\mathcal F|\mathscr F|\mathfrak F}\in\mathcal P(\mathcal P(M))$.
